% from ega0-19-fr.tex
\oldpage[0\textsubscript{IV-1}]{109}

\begin{proof}
Soit $\mathfrak m$ l'idéal maximal de $A$, de sorte que
$\mathfrak m_0=\mathfrak m/\mathfrak J$ est l'idéal maximal de $A_0$, $A$ et
$A_0$ ayant le même corps résiduel $k$. Posons
$B_{00}=B_0\otimes_{A_0}k$~; comme $B_{00}$ est une $k$-algèbre formellement
lisse (19.3.5, (iii)), c'est un anneau local régulier (19.6.5)~; en appliquant
(19.7.1.3), on voit qu'il existe une $A$-algèbre topologique $B$ qui est un
anneau local noethérien complet, un $A$-module plat, et pour laquelle on a un
$k$-isomorphisme
$u_{00}:B\otimes_A k\xrightarrow{\sim}B_{00}$. Notons qu'en vertu de (19.7.1),
$B$ est une $A$-algèbre formellement lisse, donc
$B\otimes_A A_0=B/\mathfrak JB$ est une $A_0$-algèbre formellement lisse
(19.3.5, (iii)) et un anneau local noethérien complet~; en outre, $B_0$ est un
$A_0$-module plat en vertu de l'hypothèse et de (19.7.1)~; comme on a un
$k$-isomorphisme
\[
  u_{00}:B\otimes_A k=(B\otimes_A A_0)\otimes_{A_0}k
  \xrightarrow{\sim}B_0\otimes_{A_0}k=B_{00},
\]
on déduit de (19.7.1.5), appliqué en y remplaçant $A$ par $A_0$ et
$\mathfrak J$ par $\mathfrak m_0$, qu'il existe un $A_0$-isomorphisme
$u:B\otimes_A A_0\xrightarrow{\sim}B_0$ tel que $u_{00}=u\otimes1$. Quant à
l'assertion d'unicité, notons que les idéaux $\mathfrak J^hB$ (resp.
$\mathfrak J^hB'$) sont fermés dans $B$ (resp. $B'$)
($\mathbf 0_{\mathrm I}$, 7.3.5), donc $B$ et $B'$ sont séparés et complets
pour les topologies $\mathfrak J$-préadiques (Bourbaki, \emph{Top. gén.},
chap.~III, 3\textsuperscript{e} éd., §~3, n\textsuperscript{o}~5, cor.~2 de la
prop.~9)~; on a par hypothèse un $A_0$-isomorphisme
$v_0:B\otimes_A A_0\xrightarrow{\sim}B'\otimes_A A_0$ tel que
$u'\circ v_0=u$~; comme $B$ est une $A$-algèbre formellement lisse et $B'$ un
$A$-module plat, on peut appliquer (19.7.1.5), d'où l'existence du
$A$-isomorphisme $v$ répondant à la question.
\end{proof}

\begin{remarks}[19.7.3]
\label{0.19.7.3-fr}
\begin{enumerate}
  \item[(i)] On notera que l'assertion d'unicité dans (19.7.2) est encore
  valable si l'on suppose seulement que $B$ et $B'$ sont complets \emph{pour
  les topologies $\mathfrak J$-préadiques}. Nous ignorons si l'on peut
  améliorer de même l'assertion d'existence, autrement dit si l'on peut se
  dispenser de supposer l'anneau local $B_0$ \emph{complet} (pour sa topologie
  $\mathfrak n_0$-préadique, en désignant par $\mathfrak n_0$ son idéal
  maximal) en exigeant seulement que $B$ soit complet \emph{pour la topologie
  $\mathfrak J$-préadique}. Lorsque $A_0$ est complet pour la topologie
  $\mathfrak m$-préadique, on peut voir que ce problème revient au suivant~:
  si $B_0$ est un anneau local noethérien régulier (non nécessairement
  complet) contenant le corps premier $\mathbf F_p=\mathbf Z/p\mathbf Z$,
  existe-t-il pour tout $n>1$ une $(\mathbf Z/p^n\mathbf Z)$-algèbre plate $B$
  telle que $B/pB$ soit isomorphe à $B_0$~?
  \item[(ii)] On notera qu'en général, l'isomorphisme $v$ dont l'existence est
  affirmée dans (19.7.2) \emph{n'est pas unique} (cf.~(19.8.7)).
\end{enumerate}
\end{remarks}

\subsection{Algèbres de Cohen et $p$-anneaux de Cohen; application à la
structure des anneaux locaux complets}
\label{subsection:0.19.8-fr}

Les résultats de cette section sont des applications immédiates des théorèmes
de (19.7), mais méritent d'être explicités en raison de leur importance
pratique.

\begin{definition}[19.8.1]
\label{0.19.8.1-fr}
Soient $A$, $B$ deux anneaux locaux noethériens, $\mathfrak m$ l'idéal maximal
de $A$, $k=A/\mathfrak m$ son corps résiduel,
$\varphi:A\to B$ un homomorphisme local, faisant de $B$ une $A$-algèbre. On
dit que $B$ est une \emph{$A$-algèbre de Cohen} si elle vérifie les conditions
suivantes~:
\begin{enumerate}
  \item[(i)] $B$ est un anneau \emph{complet}.
  \item[(ii)] $B$ est un $A$-module \emph{plat}.
  \item[(iii)] $B\otimes_A k$ est un \emph{corps} (autrement dit,
  $\mathfrak mB$ est l'idéal maximal de $B$) qui est une \emph{extension
  séparable} de $k$.
\end{enumerate}
\end{definition}

